% Statistical enlargement for Article V (EN).
% Main source: N19; ordered-memory antecedent: F044.
% This file does not modify the active source and is intended for a successor.
\section{Fourier duality, finite uncertainty, and ordered memory}
\label{v:n19:sec:incertidumbre-fourier}

The domain of the finite realization follows from the typed chain already
constructed in the article:
\begin{equation}
 \mathsf{APP}\longrightarrow\mathsf{TRIT}\longrightarrow\mathsf{TPK}
 \longrightarrow X_{\mathrm{enr}}
 \longrightarrow (\mathbb Z/9\mathbb Z)^3
 \xrightarrow{\,\beta_{729}\,}
 \Lambda:=(\mathbb Z/27\mathbb Z)^2
 \longrightarrow \ell^2(\Lambda).
 \label{v:n19:eq:cadena-tipificada}
\end{equation}
Here \(X_{\mathrm{enr}}\) retains route, orientation, carry, boundary,
and memory; the passage to \((\mathbb Z/9\mathbb Z)^3\) is the relevant
positional register. The map \(\beta_{729}\), defined in
\eqref{v:ant:eq:ii-area-beta729} and proved bijective in
Proposition~\ref{v:ant:prop:ii-area-beta729}, regroups the six ordered
trits into two coordinates of order twenty-seven. The positional register
and the group realization are thereby kept distinct: \(\beta_{729}\)
transports the ordered incidence, whereas the Fourier transform uses the
operations intrinsic to \((\mathbb Z/27\mathbb Z)^2\). The genealogy thus
specifies more than the common cardinality of \(729\) positions.

\subsection{Unitary characters and mutually unbiased bases}
\label{v:n19:sec:caracteres-mub}

Set \(N:=|\Lambda|=27^2=729\). For
\(x=(x_1,x_2)\) and \(k=(k_1,k_2)\) in \(\Lambda\), let
\begin{equation}
 \chi_k(x):=
 \exp\!\left(\frac{2\pi i}{27}(k_1x_1+k_2x_2)\right).
 \label{v:n19:eq:caracteres}
\end{equation}
The scalar product of the coordinates is read modulo twenty-seven.
Character orthogonality gives
\begin{equation}
 \sum_{x\in\Lambda}\chi_k(x)\overline{\chi_{\ell}(x)}
 =N\,\delta_{k,\ell}.
 \label{v:n19:eq:ortogonalidad-caracteres}
\end{equation}
Define the unitary transform \(\mathcal F_\Lambda\) by
\begin{equation}
 \widehat\psi(k):=(\mathcal F_\Lambda\psi)(k)
 =\frac1{\sqrt N}\sum_{x\in\Lambda}\chi_k(x)\psi(x),
 \qquad
 \psi(x)=\frac1{\sqrt N}\sum_{k\in\Lambda}
 \overline{\chi_k(x)}\widehat\psi(k).
 \label{v:n19:eq:fourier-inversion}
\end{equation}
The conjugated character in the second formula is indispensable under
the positive convention of the first. Equation
\eqref{v:n19:eq:ortogonalidad-caracteres} proves both inversion and
\(\|\mathcal F_\Lambda\psi\|_2=\|\psi\|_2\).

\begin{lemma}[Mutual unbiasedness of the bases]
\label{v:n19:lem:mub}
Let \(\{e_x:x\in\Lambda\}\) be the position basis and let
\(f_k(x)=N^{-1/2}\overline{\chi_k(x)}\) be the Fourier basis. Then
\begin{equation}
 |\langle e_x,f_k\rangle|=N^{-1/2}=1/27
 \qquad(x,k\in\Lambda).
 \label{v:n19:eq:mub}
\end{equation}
Thus the two bases are mutually unbiased.
\end{lemma}
\begin{proof}
The orthonormality of \(\{f_k\}\) follows from
\eqref{v:n19:eq:ortogonalidad-caracteres}; moreover,
\(\langle e_x,f_k\rangle=f_k(x)\), and every character has modulus one.
\end{proof}

\subsection{Exact support uncertainty}
\label{v:n19:sec:soporte}

For a function \(u:\Lambda\to\mathbb C\), write
\(\operatorname{supp}u=\{x:u(x)\ne0\}\).

\begin{theorem}[Support bound on the triadic torus]
\label{v:n19:thm:soporte}
Every nonzero \(\psi\in\ell^2(\Lambda)\) satisfies
\begin{equation}
 |\operatorname{supp}\psi|\,
 |\operatorname{supp}\widehat\psi|\ge N=729.
 \label{v:n19:eq:cota-soporte}
\end{equation}
\end{theorem}
\begin{proof}
Set \(S=\operatorname{supp}\psi\) and
\(T=\operatorname{supp}\widehat\psi\). The correct inversion formula in
\eqref{v:n19:eq:fourier-inversion} yields, for every \(x\in S\),
\[
 |\psi(x)|
 =\left|\frac1{\sqrt N}\sum_{k\in T}
       \overline{\chi_k(x)}\widehat\psi(k)\right|
 \le \sqrt{\frac{|T|}{N}}\,\|\widehat\psi\|_2,
\]
by Cauchy--Schwarz and \(|\chi_k(x)|=1\). Squaring and summing over
\(S\) gives
\[
 \|\psi\|_2^2
 \le\frac{|S||T|}{N}\,\|\widehat\psi\|_2^2
 =\frac{|S||T|}{N}\,\|\psi\|_2^2.
\]
The norm is positive, so cancellation proves
\eqref{v:n19:eq:cota-soporte}. This is the specialization to the realized
finite group of the Donoho--Stark support uncertainty
principle~\cite{donoho_stark_1989}.
\end{proof}

The chain \eqref{v:n19:eq:cadena-tipificada} fixes the domain before the
inequality is formulated: APP--TRIT--TPK supplies the positional register,
and \(\beta_{729}\) supplies its finite realization as the group
\(\Lambda\).

\subsection{Annihilators and classification of the saturation cases}
\label{v:n19:sec:anuladores}

For \(H\le\Lambda\), define its annihilator by
\begin{equation}
 H^\perp:=\{k\in\Lambda:\chi_k(h)=1\text{ for every }h\in H\}.
 \label{v:n19:eq:anulador}
\end{equation}
For divisors \(a,b\mid27\), consider
\begin{equation}
 H_{a,b}:=\{(x_1,x_2)\in\Lambda:
 x_1\equiv0\pmod a,\ x_2\equiv0\pmod b\}.
 \label{v:n19:eq:subgrupo-rectangular}
\end{equation}
Then
\begin{equation}
 |H_{a,b}|=\frac{27}{a}\frac{27}{b},
 \quad
 H_{a,b}^{\perp}
 =\{(k_1,k_2):k_1\equiv0\pmod{27/a},
                  k_2\equiv0\pmod{27/b}\},
 \quad |H_{a,b}^{\perp}|=ab.
 \label{v:n19:eq:anulador-explicito}
\end{equation}
Indeed, the geometric sum in each coordinate vanishes unless
\(27\mid ak_1\) and \(27\mid bk_2\). Consequently, for the normalized
indicator
\begin{equation}
 \psi_{a,b}:=|H_{a,b}|^{-1/2}\mathbf1_{H_{a,b}},
 \label{v:n19:eq:indicador-normalizado}
\end{equation}
one obtains exactly
\begin{equation}
 \widehat\psi_{a,b}(k)
 =\sqrt{\frac{|H_{a,b}|}{N}}\,
   \mathbf1_{H_{a,b}^{\perp}}(k).
 \label{v:n19:eq:fourier-indicador}
\end{equation}
Thus the transform is uniform on the annihilator and
\(|H_{a,b}|\,|H_{a,b}^{\perp}|=N\).

\begin{table}[htbp]
\centering
\caption{Four exact saturations, with distinct domains.}
\label{v:n19:tab:saturaciones}
\begin{tabular}{@{}ccccc@{}}
\toprule
\((a,b)\) & \(|H_{a,b}|\) & \(|H_{a,b}^{\perp}|\) & product
 & \((H_{\rm pos},H_{\rm F})\) \\
\midrule
\((3,3)\)  & \(81\) & \(9\)  & \(729\) & \((\log81,\log9)\) \\
\((9,9)\)  & \(9\)  & \(81\) & \(729\) & \((\log9,\log81)\) \\
\((3,9)\)  & \(27\) & \(27\) & \(729\) & \((\log27,\log27)\) \\
\((1,27)\) & \(27\) & \(27\) & \(729\) & \((\log27,\log27)\) \\
\bottomrule
\end{tabular}
\end{table}
The last two rows have the same cardinalities but neither the same subgroup
nor the same annihilator; numerical equality does not collapse their
incidence data.

\subsection{Entropic uncertainty by interpolation and differentiation}
\label{v:n19:sec:entropia}

Let \(\|\psi\|_2=1\), and define the distributions
\begin{equation}
 p_x:=|\psi(x)|^2,\qquad q_k:=|\widehat\psi(k)|^2,
 \qquad
 H_{\rm pos}:=-\sum_xp_x\log p_x,\quad
 H_{\rm F}:=-\sum_kq_k\log q_k,
 \label{v:n19:eq:entropias-distribucion}
\end{equation}
with natural logarithms and the convention \(0\log0=0\). These are the
Shannon entropies of two measurement distributions in conjugate bases.

\begin{theorem}[Entropic bound for the \(\Lambda\) realization]
\label{v:n19:thm:entropia}
Every unit vector satisfies
\begin{equation}
 H_{\rm pos}(\psi)+H_{\rm F}(\psi)
 \ge\log N=\log729=6\log3.
 \label{v:n19:eq:cota-entropica}
\end{equation}
\end{theorem}
\begin{proof}
From \(\|\mathcal F_\Lambda\|_{\ell^1\to\ell^\infty}
\le N^{-1/2}\) and unitarity,
\(\|\mathcal F_\Lambda\|_{\ell^2\to\ell^2}=1\), the Riesz--Thorin
interpolation theorem~\cite{riesz_1927,thorin_1948} gives, for
\(1<p\le2\) and \(p'=p/(p-1)\),
\begin{equation}
 \|\mathcal F_\Lambda\psi\|_{p'}
 \le N^{\frac12-\frac1p}\|\psi\|_p.
 \label{v:n19:eq:riesz-thorin}
\end{equation}
Consider
\begin{equation}
 g(p):=\log\|\mathcal F_\Lambda\psi\|_{p'}
       -\log\|\psi\|_p
       -\left(\frac12-\frac1p\right)\log N.
 \label{v:n19:eq:funcion-interpolacion}
\end{equation}
Equation \eqref{v:n19:eq:riesz-thorin} implies \(g(p)\le0\), whereas
\(g(2)=0\); hence \(g'_-(2)\ge0\).

For every unit vector \(u\) in finite dimension,
\begin{equation}
 \left.\frac{d}{dr}\log\|u\|_r\right|_{r=2}
 =\frac14\sum_j|u_j|^2\log|u_j|^2
 =-\frac14H(|u|^2).
 \label{v:n19:eq:derivada-norma}
\end{equation}
Zero coordinates are obtained by continuity, while positive terms are
differentiable. Since
\[
 \left.\frac{dp'}{dp}\right|_{p=2}=-1,
\]
left
differentiation of \eqref{v:n19:eq:funcion-interpolacion} yields
\begin{equation}
 g'_-(2)=\frac14
 \bigl(H_{\rm pos}+H_{\rm F}-\log N\bigr)\ge0,
 \label{v:n19:eq:derivada-entropica}
\end{equation}
which is equivalent to \eqref{v:n19:eq:cota-entropica}.

This derivation lies in the line of entropic uncertainty developed by
Hirschman and Beckner~\cite{hirschman_1957,beckner_1975}; for arbitrary
pairs of bases, the maximum-overlap formulation is due to
Maassen and Uffink~\cite{maassen_uffink_1988}. In the present realization,
the constant \(\log N\) follows directly by differentiating the
interpolation factor in \eqref{v:n19:eq:riesz-thorin}.
\end{proof}

For the states \(\psi_{a,b}\), Equation
\eqref{v:n19:eq:fourier-indicador} shows that \(p\) and \(q\) are uniform
on \(H_{a,b}\) and \(H_{a,b}^{\perp}\), respectively. Therefore,
\[
 H_{\rm pos}=\log|H_{a,b}|,\qquad
 H_{\rm F}=\log|H_{a,b}^{\perp}|,\qquad
 H_{\rm pos}+H_{\rm F}=\log729,
\]
which proves entropic saturation for all four cases in
Table~\ref{v:n19:tab:saturaciones} without numerical approximation.

\subsection{Typed domains of information}
\label{v:n19:sec:tipos-entropia}

The quantity \(H_{\rm pos}+H_{\rm F}\) measures the sum of two classical
distribution entropies associated with conjugate measurements.
Entanglement entropy across the boundary cut, von Neumann entropy of a
reduced state, thermodynamic entropy and temperature, and the marginal
entropy of a tritic histogram are defined on distinct objects and domains.
Every comparison among these quantities is formulated through an explicit
map specifying domain, codomain, and state. The bound
\eqref{v:n19:eq:cota-entropica} belongs entirely to
\(\ell^2(\Lambda)\).

\subsection{Ordered memory on the dodecaphasic domain}
\label{v:n19:sec:memoria-ordenada}

The F044 result acts on dodecaphasic words and characterizes the ordered
information retained by the enriched state. Its domain and operators are
independent of the Fourier realization of N19; the relation between the
two results compares marginal information with ordered information.

Let \(\tau=(\tau_1,\ldots,\tau_{12})\in\{-1,0,+1\}^{12}\), with
\(\operatorname{sgn}(0)=0\) and cyclic indices modulo twelve. Define
\begin{align}
 G_a&:=\{a,a+4,a+8\},\qquad a=1,2,3,4,\nonumber\\
 \nu_{120}(\tau)&:=\sum_{a=1}^4
   \operatorname{sgn}\!\left(\sum_{m\in G_a}\tau_m\right),
 \label{v:n19:eq:nu120}\\
 \nu_{270}(\tau)&:=\sum_{m=1}^{12}\tau_m\tau_{m+3}.
 \label{v:n19:eq:nu270}
\end{align}
The words
\begin{equation}
 \tau^{(a)}=(1,1,0,0,0,0,0,0,0,0,0,0),\qquad
 \tau^{(b)}=(1,0,0,1,0,0,0,0,0,0,0,0)
 \label{v:n19:eq:testigos-orden}
\end{equation}
have the same histogram---ten zeros and two ones---and hence the same
marginal alphabet entropy. Nevertheless,
\begin{equation}
 (\nu_{120},\nu_{270})(\tau^{(a)})=(2,0),\qquad
 (\nu_{120},\nu_{270})(\tau^{(b)})=(2,1).
 \label{v:n19:eq:testigos-lectura}
\end{equation}
The second reader detects a distance-three incidence erased by the
histogram.

For the involution \(C_M\tau:=-\operatorname{rev}(\tau)\), reversal
permutes the four classes \(G_a\), while negation reverses their sums;
the cyclic autocorrelation at lag \(+3\) equals, after reindexing, the one
at lag \(-3\). Consequently,
\begin{equation}
 \nu_{120}(C_M\tau)=-\nu_{120}(\tau),\qquad
 \nu_{270}(C_M\tau)=\nu_{270}(\tau).
 \label{v:n19:eq:paridades-lectores}
\end{equation}
The reader \(\nu_{120}\) is odd and the quadratic reader \(\nu_{270}\) is
even. These parities characterize the additional information retained by
the ordered word.

\subsection{Domain of validity and mathematical conditions}
\label{v:n19:sec:falsadores}

The conditions delimiting the two results are:
\begin{enumerate}
 \item \(\beta_{729}\) is bijective on the positional realization of
 \(\Lambda\) used in \eqref{v:n19:eq:cadena-tipificada};
 \item orthogonality \eqref{v:n19:eq:ortogonalidad-caracteres} and the
 inversion formula with the conjugated character establish unitarity;
 \item every nonzero state satisfies
 \(\lvert\operatorname{supp}\psi\rvert
 \lvert\operatorname{supp}\widehat\psi\rvert\ge729\);
 \item every unit state satisfies
 \(H_{\rm pos}+H_{\rm F}\ge\log729\);
 \item in all four cases of
 Table~\ref{v:n19:tab:saturaciones}, the transform is uniform on the
 annihilator and the cardinality product is \(729\);
 \item the witnesses in \eqref{v:n19:eq:testigos-orden} share a
 histogram, produce \eqref{v:n19:eq:testigos-lectura}, and satisfy the
 parities \eqref{v:n19:eq:paridades-lectores}.
\end{enumerate}
The first five statements belong to the typed finite realization
\(\Lambda\) and its unitary Fourier transform; the sixth belongs to the
dodecaphasic domain of ordered memory. These conditions also provide local
falsifiers: failure of any corresponding identity, inequality, or witness
value refutes the statement in its stated domain. The relation between
N19 and F044 is complementary and keeps their domains and operators
distinct.
